\section{논의}

임피던스 제어를 처음 가르칠 때 어려운 점은 수식 자체보다 수식이 표현하는 물리적 감각이다.
강성은 목표점으로 돌아가려는 힘이고, 감쇠는 빠르게 움직일수록 커지는 저항이며, 관성은 힘에 대한 가속도의 민감도를 정한다.
전기회로와 나란히 놓으면 감쇠기는 저항과, 질량은 인덕턴스와, 스프링의 순응성 $1/k$는 커패시턴스와 대응된다.
역으로 말하면 스프링 강성 $k$는 커패시턴스 $C$ 자체보다 $1/C$와 같은 위치에 놓인다.
이 세부 대응을 알고 있으면 RLC 회로와 질량-스프링-댐퍼 방정식이 왜 같은 2차 구조로 나타나는지 더 잘 보인다.

따라서 좋은 교육용 실험실에는 하나의 완성된 로봇 데모만으로는 부족하다.
오히려 작은 1차원 시스템에서 시작해 같은 개념이 어떻게 2자유도, 7자유도 시스템으로 확장되는지 보여주어야 한다.
각 단계에서 YAML 설정 파일, 로그, 그래프가 남으면 학생은 강성을 두 배로 올렸을 때 무엇이 바뀌는지, 감쇠를 줄이면 왜 진동이 생기는지, 작업공간 목표와 관절공간 목표가 무엇이 다른지를 직접 확인할 수 있다.

\subsection{학습 흐름 되짚기}

돌아보면 이 글의 흐름은 분야를 많이 넓힌 것이 아니라, 같은 질문을 조금씩 다른 물리계에서 다시 물어본 것이다.
무언가를 밀었을 때 시스템은 얼마나 움직이고, 얼마나 늦게 반응하며, 에너지를 어디에 저장하거나 잃는가?
이 질문을 전기회로에서는 전압과 전류로, 기계진동에서는 힘과 속도 또는 위치로, 로봇 제어에서는 접촉력과 끝단 운동으로 읽는다.
따라서 각 장은 따로 떨어진 배경지식이 아니라 다음 단계로 넘어가기 위한 발판이다.

\begin{table}[!htbp]
  \centering
  \caption{각 장이 다음 장으로 넘겨주는 핵심 직관}
  \begin{tabular}{p{0.18\linewidth}p{0.34\linewidth}p{0.36\linewidth}}
    \toprule
    단계 & 붙잡을 질문 & 다음 단계로 넘어가는 이유 \\
    \midrule
    임피던스의 역사 & 왜 저항 하나로는 교류 신호를 설명할 수 없었는가? & 시간에 따라 변하는 신호에는 에너지 저장과 위상 차이가 들어간다. \\
    LTI 시스템 & 복잡한 시간응답을 어떻게 입력-출력 관계로 정리할 수 있는가? & 전달함수와 주파수 응답을 쓰면 전기, 기계, 제어를 같은 문법으로 비교할 수 있다. \\
    RLC 회로 & 저항, 인덕터, 커패시터는 각각 무엇을 방해하는가? & 에너지 소산, 자기장 저장, 전기장 저장이 2차 시스템의 구조를 만든다. \\
    기계 2차계 & 같은 2차 구조가 눈에 보이는 운동에서는 어떻게 나타나는가? & 질량, 감쇠, 강성이 고유진동수, 감쇠비, 오버슈트, 정착시간을 만든다. \\
    임피던스 제어 & 로봇 끝단에 원하는 물리적 감각을 어떻게 부여하는가? & 가상 질량-스프링-댐퍼 관계를 설계해 접촉 상황의 힘과 운동을 조절한다. \\
    MuJoCo 실험실 & 이 직관을 어떻게 재현 가능한 학습 경험으로 바꾸는가? & YAML 설정 파일, 로그, 플롯을 통해 파라미터 변화와 응답 변화를 직접 비교한다. \\
    \bottomrule
  \end{tabular}
\end{table}

이 관점에서 보면 이론과 시뮬레이터는 서로 경쟁하는 설명 방식이 아니다.
이론은 어떤 파라미터가 어떤 응답을 만들지 예측하게 해주고, 시뮬레이터는 그 예측이 시간 그래프와 로봇 움직임에서 어떻게 보이는지 확인하게 해준다.
입문자에게 중요한 것은 처음부터 모든 모델을 정확히 아는 것이 아니라, 같은 개념이 표현만 바뀌어 계속 등장한다는 사실을 알아차리는 것이다.

입문자는 임피던스를 로봇 손끝에 보이지 않는 스프링과 댐퍼를 달아 두는 일로 먼저 기억해도 좋다.
목표점에서 멀어질수록 스프링 힘이 커지고, 빠르게 움직일수록 댐퍼 힘이 커진다.
이 두 힘의 합이 로봇이 버티거나 물러나는 느낌을 만든다.
여기에 관성까지 포함하면, 로봇이 외란을 받았을 때 얼마나 바로 움직이고 얼마나 묵직하게 반응하는지도 설명할 수 있다.

다음 표는 어려운 새 용어를 외우기 위한 표가 아니다.
한쪽에는 무엇이 밀고 있는가를, 다른 한쪽에는 무엇이 흐르거나 움직이는가를 놓고 보는 연습이다.
전기에서는 전압이 밀고 전류가 흐르며, 로봇에서는 외력이나 접촉력이 밀고 손끝 속도가 움직임의 흐름이 된다.
제어기가 계산한 \(\bm{f}_{\mathrm{cmd}}\)나 Lab04의 가상 벽 힘은 이 관계를 흉내 내거나 구현하기 위한 신호이지, 항상 실제 물리 포트에서 측정한 힘과 같지는 않다.

\begin{table}[!htbp]
  \centering
  \caption{전기, 기계, 로봇에서 작용 변수(effort)와 흐름 변수(flow)의 대응}
  \small
  \setlength{\tabcolsep}{4pt}
  \begin{tabularx}{\linewidth}{@{}>{\raggedright\arraybackslash}p{0.17\linewidth}>{\raggedright\arraybackslash}p{0.19\linewidth}>{\raggedright\arraybackslash}p{0.27\linewidth}>{\raggedright\arraybackslash}X@{}}
    \toprule
    관점 & 작용 변수 & 흐름/운동 상태 & 중심 질문 \\
    \midrule
    전기회로 & 전압 $v$ & 전류 $i$ & 이 주파수에서 전류가 얼마나 흐르는가? \\
    기계진동 & 힘 $f$ & 속도 $\dot{x}$ & 이 속도 운동을 만들려면 얼마나 큰 힘이 필요한가? \\
    로봇 접촉 & 접촉력/외력 & 끝단 속도 $\dot{x}$가 flow이고, 위치 $x$는 그 적분으로 생기는 오차/저장 상태 & 접촉 순간에는 속도와 감쇠를, 오래 누른 뒤에는 위치 오차와 강성을 함께 본다. \\
    \bottomrule
  \end{tabularx}
\end{table}

논의를 관통하는 메시지는 하나이다.
임피던스는 저항보다 넓은 개념이다.
어떤 시스템이 에너지를 저장하고, 소산하고, 다시 돌려주는 방식까지 포함해 힘과 운동의 관계를 설명하는 언어이다.
로봇 임피던스 제어는 이 언어를 사용해 로봇의 접촉 행동을 설계한다.

\subsection{처음 읽을 때 자주 생기는 오해}

입문자가 가장 자주 길을 잃는 지점은 용어 자체보다 층위의 변화이다.
같은 힘과 운동의 관계라도 전기회로에서는 전압과 전류로, 1차원 기계계에서는 힘과 속도 또는 위치로, 로봇에서는 끝단 오차와 접촉력으로 다시 적힌다.
아래 정리는 그 층위를 바꿔 읽을 때 흔히 생기는 착각을 모은 것이다.
표를 외우기보다는, 결과가 이상하게 보일 때 어느 구분을 놓쳤는지 되짚는 용도로 쓰면 충분하다.

\begin{table}[!htbp]
  \centering
  \caption{임피던스 학습에서 자주 생기는 개념 오해}
  \label{tab:discussion-concept-misunderstandings}
  {\renewcommand{\arraystretch}{1.12}
  \begin{tabular}{>{\raggedright\arraybackslash}p{0.27\linewidth}>{\raggedright\arraybackslash}p{0.31\linewidth}>{\raggedright\arraybackslash}p{0.32\linewidth}}
    \toprule
    오해 & 왜 헷갈리는가 & 헷갈릴 때의 기준 \\
    \midrule
    임피던스는 그냥 저항이다 & 직류 회로에서는 저항 하나로 충분해 보인다 & 임피던스는 소산뿐 아니라 에너지 저장과 위상 차이까지 포함한다 \\
    강성이 곧 임피던스이다 & 스프링에서 $f=kx$가 매우 직관적이다 & 강성은 임피던스의 스프링 성분이고, 전체 임피던스에는 감쇠와 관성도 들어간다 \\
    임피던스는 항상 힘/속도 비율만 뜻한다 & 앞에서는 기계 임피던스를 $F/V_x$로 정의하고, 뒤에서는 정적 접촉의 $F/X$와 로봇 임피던스 제어식을 함께 다룬다 & 좁은 주파수영역 기계 임피던스는 힘/속도 비율이고, 로봇 임피던스 제어는 원하는 힘-운동 관계 전체를 설계한다. 정적 접촉은 힘-위치 관계로 읽는다 \\
    순응성과 어드미턴스는 완전히 같은 말이다 & 둘 다 잘 움직이는 정도처럼 들린다 & 위치 순응성은 위치/힘, 기계 어드미턴스는 속도/힘이다 \\
    커패시터는 스프링 강성과 직접 대응된다 & 둘 다 에너지를 저장한다 & 역할은 비슷하지만 $C$는 스프링 순응성 $1/k$ 쪽과 더 가깝다 \\
    감쇠계수와 감쇠비는 같은 값이다 & 둘 다 damping이라는 이름을 가진다 & 감쇠계수는 물리 파라미터이고, 감쇠비는 $m,b,k$가 함께 만드는 무차원 비율이다 \\
    \bottomrule
  \end{tabular}}
\end{table}

표~\ref{tab:discussion-control-misunderstandings}에서는 제어와 실험 결과를 해석할 때 자주 생기는 오해를 모은다.
\begin{table}[!htbp]
  \centering
  \caption{임피던스 제어와 실험 해석에서 자주 생기는 오해}
  \label{tab:discussion-control-misunderstandings}
  {\renewcommand{\arraystretch}{1.12}
  \begin{tabular}{p{0.27\linewidth}p{0.31\linewidth}p{0.32\linewidth}}
    \toprule
    오해 & 왜 헷갈리는가 & 헷갈릴 때의 기준 \\
    \midrule
    정상상태만 보면 전체 동역학을 알 수 있다 & 최종 위치 오차는 $F/k$로 단순하다 & 접촉 순간, 오버슈트, 정착시간은 질량과 감쇠까지 봐야 한다 \\
    강성을 키우면 항상 제어가 좋아진다 & 위치 오차가 줄어드는 효과가 눈에 띈다 & 접촉력이 커지고 진동, 포화, 노이즈 민감도가 증가할 수 있다 \\
    임피던스 제어와 어드미턴스 제어는 이름만 다르다 & 둘 다 힘과 운동의 관계를 다룬다 & 구현 관점에서 임피던스는 운동을 보고 힘을 만들고, 어드미턴스는 힘을 보고 운동 명령을 만든다 \\
    작업공간 강성은 항상 그대로 실현된다 & 작업공간 강성을 설정하면 끝단 감각이 자세와 무관하다고 생각한다 & 실제 거동은 자코비안, 특이점, 관절 한계, 액추에이터 인터페이스의 영향을 받는다 \\
    시뮬레이터는 정밀 디지털 트윈이다 & Lab04 결과를 실제 Panda 접촉 성능 보증으로 읽는다 & 이 실험실은 개념을 로그와 플롯으로 확인하는 교육용 관찰 도구이다 \\
    \bottomrule
  \end{tabular}}
\end{table}

이 오해들은 대체로 같은 구분을 놓칠 때 생긴다.
정적인 관계와 동적인 관계, 힘의 크기와 부호, 위치와 속도, 물리 파라미터와 정규화된 성능 지표를 구분하지 않으면 같은 식도 서로 다른 뜻으로 읽힌다.
따라서 임피던스 제어를 공부할 때는 항상 이 식이 정상상태를 말하는지, 과도응답을 말하는지, 위치를 출력으로 보는지, 속도를 출력으로 보는지 먼저 확인하는 습관이 필요하다.

\subsection{핵심 수식과 파라미터 감각}

본문에서 반복한 수식을 한곳에 모으면 다음과 같다.
처음 실험을 해석할 때 먼저 붙잡을 기준은 다섯 가지이다.
\begin{itemize}
  \item 일정한 접촉 힘을 오래 받는 최종 위치 오차는 주로 강성이 정한다.
  \item 목표점 주변에서 흔들리고 멈추는 과정은 질량과 감쇠, 특히 감쇠비가 정한다.
  \item 첫 파라미터 숫자는 원하는 접촉력과 허용 변위에서 $k_d \approx F/\Delta x$로 잡고, 그다음 원하는 흔들림에서 $\zeta$를 고른 뒤, 내부 감쇠계수는 $d_d=2\zeta\sqrt{m_{\mathrm{nom}}k_d}$ 또는 유효질량 기반 근사로 다시 계산한다.
  \item 이때 $m_{\mathrm{nom}}$은 실제 끝단 질량의 보증값이 아니라, 교육용 데모에서 감쇠계수를 일관되게 계산하기 위한 기준 질량이다.
  \item 로그와 플롯을 볼 때는 위치 오차, 속도, 접촉력, 제어 입력이 어느 파라미터 변화와 함께 움직이는지 같이 봐야 한다.
\end{itemize}

예를 들어 손끝이 벽에 닿았을 때 \(2~\mathrm{cm}\) 정도 밀리는 동안 힘이 \(10~\mathrm{N}\)을 넘지 않게 보고 싶다면, 첫 강성값은
\[
  k_d \approx \frac{10}{0.02}=500~\mathrm{N/m}
\]
로 잡아 볼 수 있다.
그다음 그래프에서 너무 흔들리면 감쇠비를 올리고, 너무 둔하면 감쇠비를 낮춰 보며 다시 확인한다.
이 숫자는 정답이 아니라 첫 실험점을 고르는 방법이다.

이제 아래 표들은 공식 암기장이 아니라 플롯 옆에 놓고 보는 해석 메모로 생각하면 된다.
먼저 지금 보고 있는 비율이 전기 페이저 임피던스인지, 기계적 힘/속도 임피던스인지, 위치오차/힘 순응성인지, 오차속도/힘 어드미턴스인지 구분한다.
그다음 정상상태 크기는 주로 강성으로, 오버슈트와 정착 과정은 \(\zeta\), \(\omega_n\), \(\omega_d\), \(d_d\)로 나누어 읽는다.
마지막으로 로봇 실험에서는 같은 식이 \texttt{tau\_cmd\_*}, \texttt{current\_proxy\_*}, MuJoCo \texttt{data.actuator\_force}, 가상 벽 힘 같은 로그 신호로 어떻게 남았는지 대조한다.
특히 Lab04의 저장 신호는 실제 Panda 토크, 모터 전류, 외력 측정값이 아니라 교육용 프록시로 읽는다.
먼저 전기 임피던스와 기계 임피던스의 대응을 정리한다.

\begin{table}[!htbp]
  \centering
  \caption{전기와 기계 임피던스의 핵심 대응}
  \label{tab:discussion-impedance-formulas}
  {\renewcommand{\arraystretch}{1.12}
  \begin{tabular}{p{0.34\linewidth}p{0.31\linewidth}p{0.25\linewidth}}
    \toprule
    수식 & 읽는 법 & 주의할 점 \\
    \midrule
    $Z_R=R$, $Z_L=\jj\omega L$, $Z_C=1/(\jj\omega C)$ & 저항은 소산, 인덕터는 전류 변화 저항, 커패시터는 전압 변화 저항 & 정현파 정상상태의 페이저 관계이다 \\
    $Z_m(s)=F(s)/V_x(s)=ms+b+k/s$ & 힘을 속도 흐름으로 나눈 기계적 임피던스 & 분모는 위치가 아니라 속도이다 \\
    $C_e(s)=E(s)/F_{\mathrm{ext}}(s)$, $Y_e(s)=V_e(s)/F_{\mathrm{ext}}(s)=sC_e(s)$ & 오차 순응성은 위치 오차/힘, 오차 어드미턴스는 오차속도/힘이다 & 여기서 $V_e(s)=\mathcal{L}\{\dot e(t)\}$이다. 고정 목표점에서는 $V_e=V_x$라서 $Y_m=V_x/F_{\mathrm{ext}}$처럼 읽을 수 있지만, 목표점이 움직이면 $V_x=V_e+V_d$이므로 구분해야 한다 \\
    \bottomrule
  \end{tabular}}
\end{table}

다음으로 1차원 2차 시스템에서 응답 모양을 읽는 식을 정리한다.
\begin{table}[!htbp]
  \centering
  \caption{1차원 2차 시스템의 응답 감각}
  \label{tab:discussion-second-order-formulas}
  {\renewcommand{\arraystretch}{1.12}
  \begin{tabular}{p{0.34\linewidth}p{0.31\linewidth}p{0.25\linewidth}}
    \toprule
    수식 & 읽는 법 & 주의할 점 \\
    \midrule
    $m\ddot{x}+b\dot{x}+kx=f$ & 질량, 감쇠, 스프링이 힘과 운동을 나누어 담당한다 & 부호는 좌표와 힘 기준에 따라 달라질 수 있다 \\
    $\omega_n=\sqrt{k/m}$, $\omega_d=\omega_n\sqrt{1-\zeta^2}$ & 스프링이 세고 질량이 작을수록 무감쇠 기준 고유진동수가 커진다 & 피크 간격으로 보이는 부족감쇠 진동은 $\omega_d$에 가깝고, $\zeta\ge1$이면 반복 진동 주기가 없다 \\
    $\zeta=b/(2\sqrt{mk})$ & 감쇠가 임계감쇠에 비해 어느 정도인지 보는 비율 & $b$만 보고 진동 여부를 판단하면 안 된다 \\
    $d_d=2\zeta\sqrt{m_dk_d}$ 또는 $2\zeta\sqrt{m_{\mathrm{nom}}k_d}$ & 원하는 감쇠비에서 내부 감쇠계수를 계산한다 & $m_d$는 제어기가 구현하려는 목표 모델의 희망 관성이고, $m_{\mathrm{nom}}$은 실제 유효질량을 쓰지 않을 때의 교육용 기준 질량이다. 둘 다 실제 끝단 유효질량 보증값은 아니다 \\
    $U=\frac{1}{2}kx^2$, $f=-kx$ & 위치가 멀어질수록 탄성 에너지가 쌓이고 복원력이 생긴다 & 같은 힘의 크기 $|f|$라도 $k$와 $x$ 조합에 따라 저장 에너지가 달라진다 \\
    $|e_{\mathrm{ss}}|\approx |F_0|/k_d$ & 일정한 힘의 크기를 오래 받으면 강성이 최종 위치 오차의 크기를 정한다 & 과도응답에서는 $m_d$, $d_d$도 중요하다 \\
    \bottomrule
  \end{tabular}}
\end{table}

로봇 구현에서는 이상적인 목표 관계와 교육용 단순화를 구분해서 보아야 한다.
\begin{table}[!htbp]
  \centering
  \caption{로봇 임피던스 구현에서 되돌아볼 식}
  \label{tab:discussion-robot-formulas}
  {\renewcommand{\arraystretch}{1.12}
  \begin{tabular}{p{0.34\linewidth}p{0.31\linewidth}p{0.25\linewidth}}
    \toprule
    수식 & 읽는 법 & 주의할 점 \\
    \midrule
    $\begin{aligned}[t]
      &\bm{M}_d(\ddot{\bm{x}}-\ddot{\bm{x}}_d)
      +\bm{D}_d(\dot{\bm{x}}-\dot{\bm{x}}_d)\\
      &\quad+\bm{K}_d(\bm{x}-\bm{x}_d)
      =\bm{f}_{\mathrm{ext}}
    \end{aligned}$ & 원하는 질량-스프링-댐퍼 관계를 작업공간에서 쓰는 이상적 형태이다 & 실제 구현에는 동역학 보상, 토크 한계, 센서 대역폭, 안정성 조건이 필요하다 \\
    $\bm{f}_{\mathrm{cmd}}=-\bm{K}_d(\bm{x}-\bm{x}_d)-\bm{D}_d\dot{\bm{x}}$ & 끝단을 목표점에 묶인 가상 스프링-댐퍼처럼 만든다 & 고정 목표 예시이다. 움직이는 목표에서는 보통 상대속도 $\dot{\bm{x}}-\dot{\bm{x}}_d$를 쓴다 \\
    $\bm{\tau}=\bm{J}_v^{\mathsf{T}}\bm{f}$ & 병진 손끝힘을 관절토크로 바꾸는 가상일 관계이다 & 실제 힘 실현은 토크 제한, 특이점, 제어 대역폭, 모델 보상에 영향을 받는다. 6차원 렌치에는 $\bm{J}_g$를 쓴다 \\
    힘/토크 명령 가능 $\rightarrow$ 임피던스 구현, 위치 명령만 가능 + 측정/추정 외력 또는 계산된 상호작용 신호 $\rightarrow$ 어드미턴스식 외부 루프 또는 목표 위치 보정 & 같은 접촉 감각 목표도 로봇 인터페이스와 사용할 수 있는 힘 신호에 따라 구현 신호가 달라진다 & Lab04 가상 벽은 실측 외력 기반 어드미턴스가 아니라 계산된 가상 벽 신호로 목표를 후퇴시키는 교육용 근사이다 \\
    \bottomrule
  \end{tabular}}
\end{table}

이 표의 첫 번째 식은 이상적인 작업공간 임피던스 관계를 보여 주고, 두 번째 식은 튜토리얼과 실험실에서 자주 쓰는 스프링-댐퍼 근사이다.
따라서 Lab04 결과는 실제 로봇 성능의 판정표라기보다, 강성, 감쇠, 목표점, 접촉력이 시간응답과 플롯에 어떤 흔적으로 남는지 읽어 보는 교육용 눈금자로 삼으면 된다.

\subsection{목적에 따라 달라지는 파라미터 선택}

강성, 감쇠, 관성은 단순히 크게 또는 작게 만들면 되는 값이 아니다.
강성을 키우면 위치 오차는 줄어들지만 접촉력이 커질 수 있고, 강성을 낮추면 접촉력 변화가 완만해질 수 있지만 원하는 위치에서 더 많이 밀릴 수 있다.
낮은 강성이 곧바로 안전을 뜻하는 것은 아니며, 실제 안전성은 속도 제한, 힘 제한, 감쇠, 환경 강성, 제어 대역폭까지 함께 보아야 한다.
감쇠를 늘리면 진동과 오버슈트가 줄어들지만, 너무 크면 움직임이 둔해지거나 위치 제어기가 더 큰 노력을 요구할 수 있다.
관성은 빠른 외란과 충격에 대한 반응을 바꾸지만, 실제 로봇에서는 모델 오차와 구동기 한계 때문에 사용자가 직접 조절하기 어려운 경우가 많다.

임피던스 제어를 잘 이해하려면 식을 외우는 데서 멈추면 안 된다.
목표가 빠른 도달인지, 작은 접촉력인지, 진동 없는 정착인지, 사람에게 안전한 순응성인지에 따라 어떤 파라미터를 바꾸어야 하는지 말할 수 있어야 한다.
이 튜토리얼과 실험실은 그 파라미터 선택 기준을 익히기 위한 장치이다.

처음 튜닝할 때 가장 실행 가능한 순서는 네 단계이다.
먼저 접촉 방향을 고른다.
그다음 원하는 힘과 허용 변위로 $k_d \approx F/\Delta x$를 잡는다.
이후 빠르지만 약간 흔들려도 되는지, 오버슈트를 줄일지에 따라 $\zeta$를 고르고, 내부에서는 명목 질량 또는 유효질량으로 $d_d$를 다시 계산한다.
마지막으로 토크나 힘 명령이 가능한지, 위치 명령만 가능한지, 그리고 측정/추정 외력이나 Lab04처럼 계산된 상호작용 신호를 사용할 수 있는지에 따라 임피던스 구현, 어드미턴스식 외부 루프, 또는 목표 위치 보정 구현을 고른다.

이 글을 끝까지 읽은 뒤의 목표는 모든 수식을 외우는 것이 아니다.
독자는 먼저 이 상황에서는 무엇이 원인이고 무엇이 반응인지, 에너지는 어디에 저장되고 어디에서 사라지는지, 강성이나 감쇠를 바꾸면 그래프의 어느 부분이 달라지는지를 스스로 물을 수 있으면 된다.
그다음 MuJoCo 실험실에서는 같은 질문을 설정 파일의 파라미터, 로그의 신호, 플롯의 응답 모양에 차례로 연결해 확인한다.
이 튜토리얼은 개념을 외우게 하려는 목차라기보다, 예측하고 실행하고 그래프로 확인하게 하는 흐름이다.

